% 原章第80页结论7.27—第81页7.3.3.4标题前。
\Needspace{5\baselineskip}
\noindent\textbf{结论 7.27\quad 高斯过程回归}

\noindent\textbf{步骤1。}定义高斯核矩阵 $K(X,X)$，计算噪声方差初值：
\[
(\widehat\sigma_n^2)^{(0)}=\frac1{n-q}\sum_{i=1}^{n}(y_i-\overline y)^2,
\qquad \overline y=\frac1n\sum_{i=1}^{n}y_i.
\]
令 $t=0$。

\noindent\textbf{步骤2。}估计回归系数：
\[
\widehat\alpha=\bigl(K(X,X)+(\widehat\sigma_n^2)^{(t)}I\bigr)^{-1}y.
\]

\noindent\textbf{步骤3。}给定 $x_{new}$，计算预测值 $\overline y_{new}=k(x_{new},X)^T\widehat\alpha$ 和噪声 $e_{new}=y_{new}-\overline y_{new}$。

\noindent\textbf{步骤4。}计算预测值估计量的协方差矩阵，以及测试数据 $y=y_{new}$ 的预测分布：
\[
V(\widehat y_{new})=K(X_{new},X_{new})-K(X_{new},X)K(X,X)^{-1}K(X,X_{new}),
\]
\[
\begin{aligned}
p(y\mid X)={}&\frac1{(2\pi)^{n/2}|K(X,X)+\sigma_n^2I|^{1/2}}\\
&\quad\times\exp\left\{-\frac12y^T[K(X,X)+\sigma_n^2I]^{-1}y\right\}.
\end{aligned}
\]

\noindent\textbf{步骤5。}计算测试数据集的噪声方差：
\[
(\widehat\sigma_n^2)^{(t)}=\frac1{m-q}\sum_{i=1}^{m}
\bigl(y_{new_i}-k(x_{new_i},X)^T\widehat\alpha)^2\bigr).
\]

\noindent\textbf{步骤6。}检查收敛。令 $\varepsilon$ 为预设误差。如果
\[
\bigl|(\widehat\sigma_n^2)^{(t)}-(\widehat\sigma_n^2)^{(t-1)}\bigr|<\varepsilon,
\]
则停止；否则返回步骤2。

如果用 Cholesky 分解计算矩阵的逆，则结论7.27可化为结论7.28（Rasmussen and Williams 2006）。

\Needspace{5\baselineskip}
\noindent\textbf{结论 7.28\quad 使用 Cholesky 分解的高斯过程回归}

\noindent\textbf{步骤1。}定义高斯核矩阵 $K(X,X)$，计算噪声方差初值：
\[
(\widehat\sigma_n^2)^{(0)}=\frac1{n-q}\sum_{i=1}^{n}(y_i-\overline y)^2,
\qquad \overline y=\frac1n\sum_{i=1}^{n}y_i.
\]
令 $t=0$。

\noindent\textbf{步骤2。}估计回归系数。

进行 Cholesky 分解：
\[
K(X,X)+(\widehat\sigma_n^2)^{(t)}I=LL^T.
\]
解方程 $Lz=y$，求得 $z$；解方程 $L^T\alpha=z$，求得 $\alpha$。

\noindent\textbf{步骤3。}给定 $x_{new}$，计算预测值 $\overline y_{new}=k(x_{new},X)^T\widehat\alpha$ 和噪声 $e_{new}=y_{new}-\overline y_{new}$。

\noindent\textbf{步骤4。}计算预测值估计量的协方差矩阵与边际似然。

解方程 $L\gamma=K(X,X_{new})$，求得 $\gamma$。计算方差：
\[
V(\widehat y_{new})=K(X_{new},X_{new})-\gamma^T\gamma.
\]
计算带噪声测试数据 $y=y_{new}$ 的预测分布：
\[
p(y\mid X)=\frac1{(2\pi)^{n/2}|L|}\exp\left\{-\frac12y^T\widehat\alpha\right\}.
\]

\noindent\textbf{步骤5。}计算测试数据集的噪声方差：
\[
(\widehat\sigma_n^2)^{(t)}=\frac1{m-q}\sum_{i=1}^{m}
\bigl(y_{new_i}-k(x_{new_i},X)^T\widehat\alpha\bigr)^2.
\]

\noindent\textbf{步骤6。}检查收敛。令 $\varepsilon$ 为预设误差。如果
\[
\bigl|(\widehat\sigma_n^2)^{(t)}-(\widehat\sigma_n^2)^{(t-1)}\bigr|<\varepsilon,
\]
则停止；否则返回步骤2。
